Les présentes notes constituent exclusivement une synthèse du formalisme présenté dans
le cadre du cycle de cours « Physique des particules : introduction au Modèle standard »
[faisant partie de la mineure « Physique nucléaire et des particules »] dispensé en 1re
année de master de physique générale à l’Université Paris-Saclay. Les spinors et les matrices
de Dirac y sont définis et décrits en détail. Leurs transformations de Lorentz sont abordées
afin de construire des lagrangiens invariants sous Lorentz. L’équation de Dirac est étudiée
et résolue explicitement. Les opérateurs d’hélicité et de conjugaison de charge sont
présentés. Le mécanisme de Higgs générant des masses pour les fermions est expliqué,
ainsi que le théorème de Goldstone.
